\documentclass[11pt]{article}
\usepackage[utf8]{inputenc}
\usepackage[margin=1in]{geometry}
\usepackage{graphicx} % Required for inserting images
\usepackage{amsfonts}
\usepackage{amsmath}
\usepackage{dirtytalk}
\usepackage{algorithm}
\usepackage{mathtools}
\usepackage{xcolor}
\usepackage{amssymb}
\usepackage[title]{appendix}
\usepackage[noend]{algpseudocode}
\usepackage{amsthm}
\usepackage{bigints}
\usepackage{thm-restate}
\usepackage[]{hyperref}
\usepackage{cleveref}
\usepackage{appendix}
\usepackage{makecell}
\usepackage{bigints}
\usepackage{mathtools}
\usepackage{subfigure}
\usepackage{booktabs}
\usepackage{bbding}
\usepackage{wrapfig}
\usepackage{bbm}
\usepackage{authblk}
\usepackage{cleveref}

\declaretheorem[name=Theorem,numberwithin=section]{thm}
\declaretheorem[name=Lemma,numberwithin=section]{lemma}
\declaretheorem[name=Assumption,numberwithin=section]{asmpt}
\declaretheorem[name=Proposition,numberwithin=section]{prop}
\declaretheorem[name=Corollary,numberwithin=section]{corollary}
\declaretheorem[name=Remark,numberwithin=section]{rmk}
\declaretheorem[name=Definition,numberwithin=section]{defn}
\declaretheorem[name=Claim,sibling=thm]{claim}


\usepackage[numbers]{natbib}
\bibliographystyle{unsrtnat}


\newcommand{\expnumber}[2]{{#1}\mathrm{e}{#2}}
\newcommand{\man}{\mathcal{M}}
\newcommand{\norm}[1]{\left\lVert#1\right\rVert}


\title{Locally Linear Single Index Regression}

\author{Tristan Luca Saidi}
\author{Arun Kuchibhotla}

\affil[1]{Department of Statistics and Data Science, Carnegie Mellon University}
\begin{document}


\maketitle

\section{Introduction}
\subsection{Setting}
We consider the regression setting where we are given samples $\{(X_1, Y_1), \dots, (X_d, Y_n)\}$ from some joint distribution over $\mathbb{R}^d\times \mathbb{R}$. In the single-index setting, we know that the conditional expectation of $Y$ given $X$ has the following compositional structure,
\[\mathbb{E}[Y | X = x] = f(\Pi(x))\]
where $\Pi:\mathbb{R}^d \rightarrow \Gamma$ and $f: \Gamma\rightarrow\mathbb{R}$. The classical single-index setting assumes that $\Gamma = \mathbb{R}$, $\Pi(x) = b^{T}x$ for some $b \in \mathbb{R}^d$, and thus can be thought of as a projection onto a line. In this work, we consider the setting where $\Pi$ is a closest-point projection onto a compact regular $C^3$ curve $\gamma$. We choose to represent $\gamma$ as a function $\gamma:I \rightarrow \Gamma \subset \mathbb{R}^d$ with $I = [0,\text{len}_{\gamma}]$, where $\gamma$ is parameterized by arclength $t$. We assume that $\gamma$ is injective, which ensures that $\gamma^{-1}:\Gamma\rightarrow I$ is well defined and rules out any self-intersections. We define the closest-point projection map as
\begin{equation}\label{def: projection}
\Pi(x) = \arg\inf_{t\in I}\|x - \gamma(t)\|_2.
\end{equation}
\begin{restatable}[Infimum is attained]{rmk}{} \label{remark: infimum}
    The infimum in \cref{def: projection} is attained by at least one point in the set $I$ for every $x \in \mathbb{R}^d$.
\end{restatable}
\noindent \textit{Proof.} Since the set $\Gamma$ is compact and the Euclidean distance $\|x-a\|_2$ is continuous, the infimum must be attained on some element in $\Gamma$. The result follows from the injectivity of $\gamma$. 

\qed{}

\subsection{Geometry of curves in $\mathbb{R}^d$}
Before presenting our estimator and its properties, we will remark on relevant results from differential geometry and geometric measure theory. 


\begin{restatable}[Medial axis, \cite{federer1959curvature}]{defn}{}
    Given a closed subset $A \subset \mathbb{R}^d$, the medial axis $M(A)$ is the subset of $\mathbb{R}^d$ consisting of points with at least two nearest neighbors on $A$. Formally,
    \begin{equation}
    M(A) = \{x \in \mathbb{R}^d : |\{y \in A : \|x-y\|_2 = d(x,A)\}| > 1\}. \label{eq: medial axis definition}
    \end{equation}
\end{restatable}

\begin{restatable}[Reach, \cite{federer1959curvature}]{defn}{}
    The reach of a closed subset $A \subset \mathbb{R}^d$ is defined as 
    \begin{equation}
        \tau_A = \inf_{p\in A} d(p,M(A)) = \inf_{z\in M(A)}d(z,A). \label{eq: reach definition}
    \end{equation}
\end{restatable}


\begin{restatable}[Boundedness of $\gamma^{(k)}$]{rmk}{} \label{remark: deriv bound}
    The curve $\gamma$ has bounded first, second and third derivatives. Namely, for any $\gamma \in C^3(I)$ and $k \in \{1,2,3\}$ we have 
    \[\sup_{t\in I}\|\gamma^{(k)}(t)\|_2 \leq M_k\]
    for some $|M_k| < \infty$.
\end{restatable}
\noindent \textit{Proof.} Since $\gamma \in C^3(I)$, $\gamma^{(k)}$ is continuous for $k \leq 3$. Since $I$ is a compact set, by the Extreme Value Theorem $\gamma^{(k)}$ must be bounded on $I$. 

\qed{}


\noindent We will now discuss some properties of the closest point projection map $\Pi$. 


\begin{restatable}[Differentiability of $\Pi$]{thm}{} \label{thm: differentiability of pi}
    There exists an \textit{open} set $R$ of full $d$-dimensional Lebesgue measure such that the following conditions hold for all $x \in R$:
    \begin{enumerate}
        \item $s = \Pi(x)$ is unique and well defined
        \item and both $\nabla\Pi(x)$, $\nabla^2\Pi(x)$ exist and can be expressed as $\nabla\Pi(x)  = -\gamma'(s)/g(x)$ and 
        \begin{equation}
            \nabla^2\Pi(x) = \frac{1}{g^2(x)}\left(\gamma''(s)\gamma'(s)^T + \gamma'(s)\gamma''(s)^T\right)- \frac{\langle\gamma^{(3)}(s), x - \gamma(s)\rangle}{g^3(x)}\gamma'(s)\gamma'(s)^T.
        \end{equation}
        where $g(x) = -1 + \langle\gamma''(s), x-\gamma(s)\rangle$. 
    \end{enumerate}
\end{restatable}
\noindent \textit{Proof.} Let $B$ and $\overline{\text{Sing}(\Gamma)}$ be defined as in \cref{eq: ift degeneracy set} and \cref{eq: singular set} respectively. By \Cref{lemma: IFT degen} and \Cref{lemma: singular set measure} we know that both sets have Lebesgue measure $0$. By \Cref{lemma: IFT degen closed} we also know that $B$ is closed. It follows that $B \, \cup \overline{\text{Sing}(\Gamma)}$ has measure $0$ and is also closed. It follows that 
\begin{equation}
    R := \left(B \, \cup \overline{\text{Sing}(\Gamma)}\right)^c \label{eq: good set}
\end{equation}
is open with full measure. We also know that for any $x \in R$
\begin{enumerate}
    \item $x \notin \overline{\text{Sing}(\Gamma)}$, which means $\Pi(x)$ is differentiable at $x$ and $\Pi(x)$ therefore exists and is unique. 
    \item $x \notin B$, which allows us to apply \Cref{prop: proj gradient} and \Cref{prop: proj hessian} to get the expressions for $\nabla\Pi(x)$ and $\nabla^2\Pi(x)$ in the theorem statement.
\end{enumerate}
\qed{}

% \begin{restatable}[Lebesgue almost everywhere uniqueness of $\Pi$, \cite{erdos1945some}]{corollary}{} \label{remark: uniqueness}
%     The closest point projection $\Pi$ onto the closed regular $C^3$ curve $\gamma$ is Lebesgue almost everywhere unique. 
% \end{restatable}


\subsection{Distribution of covariates}
Previous works utilize the Linear Conditional Expectation (LCE) property, where we assume that the following condition on the distribution of the covariates holds: for any $b \in \mathbb{R}^p$ and $B \in \mathbb{R}^{d\times p}$, then there exists a $c \in \mathbb{R}^d$ such that the conditional expectation $\mathbb{E}[b^{T}X \, | B^{T}X] = B^{T}Xc$. In this work we will only assume that the support of $X$ is contained in $\mathcal{B}(0, r)$ for some $r > 0$. We will also suppose that $\Gamma  = \text{Im}(\gamma) \subset \mathcal{B}(0, r)$ as well. 

\begin{restatable}[Elliptically Symmetric Distributions]{defn}{}\label{def: elliptical symmetry}
Let $\mu \in \mathbb{R}^d$ and let $\Lambda \in \mathbb{R}^{d\times d}$ with $\Lambda \succeq 0$. If $X$ has density
\begin{equation}
     f(X) = |\Lambda|^{-1/2}g\left((X-\mu)^{T}\Lambda^{-1}(X - \mu)\right)
\end{equation}
and $g:[0, \infty) \rightarrow[0, \infty)$ does not depend on $\mu$, $\Lambda$, then the distribution of $X$ is said be elliptically symmetric. 
\end{restatable}

\begin{restatable}[Spherically Symmetric Distributions]{defn}{}\label{def: spherical symmetry}
If $X\in \mathbb{R}^d$ has density satisfying \Cref{def: elliptical symmetry} with $\Lambda = I_d$, then the distribution of $X$ is said to be spherically symmetric. Equivalently, a random variable $X$ has a spherically symmetric distribution about its mean $\mu$ if for every orthogonal matrix $Q$, we have $Q(X - \mu) \overset{d}{=} X - \mu$ (where $\overset{d}{=}$ denotes equivalence in distribution). 
\end{restatable}

\subsection{Recovering the underlying curve}
For a given $x_0 \in \mathbb{R}^d$ will construct an estimator $\widehat{\theta}_{x_0} \approx \gamma'(\Pi(x_0))$ using the following procedure. Let $\epsilon > 0$ and define
\begin{align} 
\widehat{\theta}_{x_0,\epsilon}, \widehat{\alpha}_{x_0, \epsilon} &= \arg\min_{\theta \in \mathbb{R}^d,\, \alpha \in \mathbb{R}} \sum_{X_i \in \mathcal{B}(x_0, \epsilon)}\frac{p^*(X_i - x_0)}{\hat{p}(X_i)}\left(Y_i - (X_i - x_0)^{T}\theta - \alpha\right)^2 \\
&= \arg\min_{\theta \in \mathbb{R}^d,\, \alpha \in \mathbb{R}} \sum_{i = 1}^{n}\mathbbm{1}\left\{X_i \in \mathcal{B}(x_0, \epsilon)\right\}\frac{p^*(X_i - x_0)}{\hat{p}(X_i)}\left(Y_i - (X_i - x_0)^{T}\theta - \alpha\right)^2
\end{align}
where $\hat{p}$ is some consistent estimate of $p$, the density of $X$. We also define $p^*$ to be any zero-mean spherically symmetric distribution with covariance $\sigma^2I_d$ for some $\sigma > 0$. By a slight abuse of notation, we'll redefine each $X_i$ to have $x_0$ subtracted off, and the optimization becomes 
\begin{align} 
\widehat{\theta}_{x_0,\epsilon}, \widehat{\alpha}_{x_0, \epsilon} &= \arg\min_{\theta \in \mathbb{R}^d,\, \alpha \in \mathbb{R}} \sum_{i = 1}^{n}\underbrace{\mathbbm{1}\left\{X_i \in \mathcal{B}(0, \epsilon)\right\}\frac{p^*(X_i)}{\hat{p}_0(X_i)}}_{\hat{w}(X_i)}\left(Y_i - X_i^{T}\theta - \alpha\right)^2
\end{align}
where $\hat{p}_0$ is now a consistent estimate of the density of $X - x_0$. Define the weighted means of $X_i$ and $Y_i$ as 
\begin{equation}
    \overline{X}_{\hat{w}} = \frac{\sum_{i = 1}^n \hat{w}(X_i)X_i}{\sum_{i = 1}^n \hat{w}(X_i)} \qquad \text{and} \qquad \overline{Y}_{\hat{w}} = \frac{\sum_{i = 1}^n \hat{w}(X_i)f(\Pi(X_i+x_0))}{\sum_{i = 1}^n \hat{w}(X_i)}.
\end{equation}
Through a straightforward calculation one would find that
\begin{align}
    \widehat{\theta}_{x_0,\epsilon} &= \left(\sum_{i = 1}^n\hat{w}(X_i)\tilde{X_i}\tilde{X_i}^T\right)^{-1}\left(\sum_{i = 1}^n\hat{w}(X_i)\tilde{X_i}\tilde{Y_i}\right) \label{eq: sample est closed form}
\end{align}
where $\tilde{X_i} = X_i - \overline{X}_{\hat{w}}$ and $\tilde{Y_i} = Y_i - \overline{Y}_{\hat{w}}$. A similar set of calculations leads to the population quantity. We will again abuse notation and assume that $X$ has had $x_0$ subtracted off. The population optimization problem then looks like
\begin{align}
\theta^*_{x_0, \epsilon}, \alpha^*_{x_0, \epsilon} &= \arg\min_{\theta, \alpha} \mathbb{E}_{X\sim p_0} \Big[\mathbbm{1}\Big\{X \in \mathcal{B}(0, \epsilon)\Big\}\frac{p^*(X)}{p_0(X)}\Big(f(\Pi(X + x_0)) - X^{T}\theta - \alpha \Big)^2\Big] \label{eq: population estimator}
\end{align}
where $p_0$ is the density of the random variable $X$. Solving for $\theta^*_{x_0, \epsilon}$ yields
\begin{align}
    \theta^*_{x_0, \epsilon} = \mathbb{E}_{X \sim p^*}\left[\mathbbm{1}\{X \in \mathcal{B}(0, \epsilon)\}XX^T\right]^{-1}\mathbb{E}_{X \sim p^*}\left[\mathbbm{1}\{X \in \mathcal{B}(0, \epsilon)\}X\,\overline{Y}\right] \label{eq: population est closed form}
\end{align}
where $\overline{Y} = f(\Pi(X+x_0)) - \mathbb{E}_{X \sim p^*}\left[\mathbbm{1}\{X \in \mathcal{B}(0, \epsilon)\}f(\Pi(X - x_0))\right]$. We will show that $\widehat{\theta}_{x_0,\epsilon}$ is a consistent estimate of $\gamma'(\Pi(x_0))$.



\begin{restatable}[Population quantity proportionality]{thm}{}
    For some scalar $C \in \mathbb{R}$ and Lebesgue almost every $x_0 \in \mathcal{B}(0, r)$ there exists an $\epsilon > 0$ such that 
    \[\|\theta^*_{x_0, \epsilon} - C\gamma'(\Pi(x_0))\|_2 \lesssim (M_2+rM_3)L_f\epsilon\]
    with $\theta^*_{x_0, \epsilon}$ defined in \cref{eq: population estimator}.
\end{restatable}

\noindent \textit{Proof.} By \Cref{thm: differentiability of pi} we have that for Lebesgue almost every $x_0$, $\Pi(x_0)$ exists, is unique and $\nabla\Pi(x_0) \propto \gamma'(\Pi(x_0))$. Thus it suffices to show that $\theta^*_{x_0, \epsilon} \rightarrow \nabla\Pi(x_0)$. To do so, we'll define 
\[g(x) := f\Big(\Pi(x_0) + \nabla\Pi(x_0)^{T}(x-x_0)\Big)\]
and 
\[\xi(x) := f(\Pi(x)) - g(x).\] Note that \Cref{lemma: proj approx} implies that there exists an $\epsilon > 0$ such that for all $x \in \mathcal{B}(x_0, \epsilon)$ we have $\xi(x) \lesssim (M_2 + rM_3)L_f\epsilon^2$. Let  
\begin{align}
    \tilde{\theta}^*_{x_0, \epsilon}, \tilde{\alpha}_{x_0}^* &= \arg\min_{\theta, \alpha}\mathbb{E}_{X\sim p_0}\Big[\underbrace{\mathbbm{1}\Big\{X \in \mathcal{B}(0, \epsilon)\Big\}\frac{p^*(X)}{p_{0}(X)}}_{:=w(X)}\Big(g(X+x_0) - X^{T}\theta - \alpha\Big)^2\Big]. \label{eq: perturbed pop estimator}
\end{align}
By \Cref{lemma: surrogate problem} we know that for some $C \in \mathbb{R}$ we have $\tilde{\theta}^*_{x_0, \epsilon}  = C\cdot \nabla\Pi(x_0)$, and by \Cref{lemma: proj approx}, we know that
\begin{align*}
    \left|\xi(X+x_0)\right| &= \big|f(\Pi(\underbrace{X+x_0}_{Z})) - f\left(\Pi(x_0) + \nabla\Pi(x_0)^TX\right)\big| \\
    &= \left|f(\Pi(Z)) - f\left(\Pi(x_0) + \nabla\Pi(x_0)^T(Z-x_0)\right)\right|\\
    &\lesssim (M_2 + rM_3)L_f\epsilon^2.
\end{align*}
Observe that $\mathbb{E}_{X\sim p_{0}}[w(X)XX^{T}]$ is necessarily positive definite, as it can be written as an expectation over $p^*$, which has no collinearity structure. Thus, \Cref{lemma: perturbation} allows us to say 
\begin{align*} 
\|\theta^*_{x_0, \epsilon} - \tilde{\theta}_{x_0, \epsilon}^*\|_2 &= \left\|\left\{\mathbb{E}_{X\sim p_{0}}[w(X)XX^{T}]\right\}^{-1}\mathbb{E}_{X\sim p_{0}}[w(X)X\xi(X+x_0)]\right\|_2 \\
&\leq \left\|\left\{\mathbb{E}_{X\sim p_{0}}[w(X)XX^{T}]\right\}^{-1}\mathbb{E}_{X\sim p_{0}}[w(X)X]\right\|_2(M_2 + rM_3)L_f\epsilon^2.
\end{align*}
Now we'll apply \Cref{lemma: pnorm bound} to obtain,
\begin{align}
    \|\theta^*_{x_0, \epsilon} - \tilde{\theta}_{x_0, \epsilon}^*\|_2 &\lesssim (M_2 + rM_3)L_f\epsilon.
\end{align}
\qed{}

\begin{restatable}[Central Limit for $\hat{\theta}_{x_0, \epsilon}$]{thm}{}
    Under \textcolor{red}{some regularity conditions,} we have asymptotic normality. In particular
    \[\sqrt{n}(\hat{\theta}_{x_0, \epsilon} - \theta^*_{x_0, \epsilon}) \overset{d}{\rightarrow} N(0, V).\]  
\end{restatable}
\noindent \textit{Proof.} This proof will follow closely the argument from \citet{Newey_Ruud_2005}. Let's define the matrices
\begin{equation}
    Q = \mathbb{E}_{X \sim p^*}\left[\mathbbm{1}\{X \in \mathcal{B}(0, \epsilon)\}XX^T\right] \qquad \text{and} \qquad \hat{Q} = \frac{1}{n}\sum_{i = 1}^n\hat{w}(X_i)\tilde{X_i}\tilde{X_i}^T.
\end{equation}
Taking the root-$n$ scaled difference between $\hat{\theta}_{x_0, \epsilon}$ and $\theta^*_{x_0, \epsilon}$ yields
\begin{align*}
    \sqrt{n}(\hat{\theta}_{x_0, \epsilon} -\theta^*_{x_0, \epsilon}) &= \sqrt{n}\left(n\hat{Q}\right)^{-1}\left(\sum_{i = 1}^n\hat{w}(X_i)X\right) - \sqrt{n}\,\theta_{x_0, \epsilon}^*\\
    &= \frac{1}{\sqrt{n}}\hat{Q}^{-1}\left(\sum_{i = 1}^n\hat{w}(X_i)\tilde{X_i}\tilde{Y_i}\right) - \sqrt{n}\hat{Q}^{-1}\hat{Q}\theta_{x_0, \epsilon}^* \\
    &= \hat{Q}^{-1}\underbrace{\frac{1}{\sqrt{n}}\sum_{i=1}^n \hat{w}(X_i)\tilde{X}_i\left(\tilde{Y}_i - \tilde{X}_i^T\theta^*_{x_0, \epsilon}\right)}_{\hat{m}}.
\end{align*}
Under appropriate regularity conditions (\textcolor{red}{what are they??}) $\hat{Q}^{-1} \overset{p}{\rightarrow} Q^{-1}$. By \textcolor{red}{need this result} we have $\hat{m} \overset{d}{\rightarrow} N(0, \Lambda)$ for some $\Lambda$ that we will derive below.  Thus, by Slutsky's theorem the asymptotic variance of $\sqrt{n}(\hat{\theta}_{x_0, \epsilon} -\theta^*_{x_0, \epsilon})$ becomes $Q^{-1}\Lambda Q^{-1}$. Deriving an expression for $\Lambda$ requires an analysis of $\hat{p}_0$, the nonparametric density estimate of $p_0$. \citet{newey1994kernel} showed that under some regularity assumptions (\textcolor{red}{what?}) on the density estimator $\hat{p}$ of density $p_0$, we have the following: for some function $m(Z, p)$ of a random variable $Z$ and scalar $p$,
\begin{align}
    \sum_{i = 1}^n\frac{m(Z_i, \hat{p}(x_i))}{\sqrt{n}} = \sum_{i = 1}^n \left[ \frac{m(Z_i, p_0(x_i)) + D(x_i)p_0(x_i) - \mathbb{E}[D(x)p_0(x)]}{\sqrt{n}}\right] + o_p(1)
\end{align}
where 
\begin{equation}
    D(x) = \mathbb{E}\left[ \frac{\partial m(Z,p) }{\partial p}\Bigg|_{p = p_0(x)}\Bigg| \, x\right].
\end{equation}
Now we'll instantiate this result with 
\[m(Z, p) = \frac{\mathbbm{1}\{X \in \mathcal{B}(0, \epsilon)\}p^*(X)X(Y-X^T\theta^*_{x_0, \epsilon})}{p}.\]
It follows that
\begin{align*}
    D(X) &= \mathbb{E}\left[ \frac{\partial m(Z,p) }{\partial p}\Bigg|_{p = p_0(X)}\Bigg| \, X\right] \\
    &= -\mathbb{E}\left[ \frac{\mathbbm{1}\{X \in \mathcal{B}(0, \epsilon)\}p^*(X)X(Y-X^T\theta^*_{x_0, \epsilon})}{p_0(X)^2} \Bigg| \, X\right] \\
    &= -p_0(X)^{-1}w(X)X(Y - X^T\theta^*_{x_0, \epsilon}).
\end{align*}
Plugging this in yields
\begin{align}
    \sum_{i = 1}^n\frac{\hat{w}(X_i)\tilde{X}_i\tilde{u}_i}{\sqrt{n}} = n^{-1/2}\sum_{i = 1}^n \left[ w(X_i)\tilde{X}_i\tilde{u}_i + D(x_i)p_0(x_i) - \mathbb{E}_{p_0}[w(X)X(Y - X^T\theta^*_{x_0, \epsilon})\right] + o_p(1)
\end{align}

\begin{restatable}{lemma}{}
    The term
    \[\frac{1}{\sqrt{n}}\sum_{i=1}^n \hat{w}(X_i)\tilde{X}_i\left(\tilde{Y}_i - \tilde{X}_i^T\theta^*_{x_0, \epsilon}\right)\]
    has a limiting normal distribution $N(0, \Lambda)$ for some $\Lambda \succeq 0$. 
\end{restatable}
\noindent \textit{Proof.} Recall the definitions $\tilde{X}_i = X_i - \overline{X}_{\hat{w}}$ and $\tilde{Y}_i = Y_i - \overline{Y}_{\hat{w}}$
where
\begin{equation}
    \overline{X}_{\hat{w}} = \frac{\sum_{i = 1}^n \hat{w}(X_i)X_i}{\sum_{i = 1}^n \hat{w}(X_i)} \qquad \text{and} \qquad \overline{Y}_{\hat{w}} = \frac{\sum_{i = 1}^n \hat{w}(X_i)f(\Pi(X_i + x_0))}{\sum_{i = 1}^n \hat{w}(X_i)}.
\end{equation}
Observe that $\frac{1}{n}\sum_{i = 1}^n\hat{w}(X_i)X_i \overset{p}{\rightarrow} \mathbb{E}_{p^*}[\mathbbm{1}\{X \in \mathcal{B}(0, \epsilon)\}X]$, $\frac{1}{n}\sum_{i = 1}^n\hat{w}(X_i)f(\Pi(X_i + x_0)) \overset{p}{\rightarrow} \mathbb{E}_{p^*}[\mathbbm{1}\{X \in \mathcal{B}(0, \epsilon)\}f(\Pi(X+x_0))]$ and $\frac{1}{n}\sum_{i = 1}^n\hat{w}(X_i) \overset{p}{\rightarrow} \mathbb{E}_{p^*}[\mathbbm{1}\{X \in \mathcal{B}(0, \epsilon)\}] = \mathbb{P}_{p^*}[X \in \mathcal{B}(0, \epsilon)]$. For $\epsilon > 0$ this quantity is nonzero, which allows us to apply Slutsky's theorem to say
\[\overline{X}_{\hat{w}} \overset{p}{\rightarrow} \frac{\mathbb{E}_{p^*}[\mathbbm{1}\{X \in \mathcal{B}(0, \epsilon)\}X]}{\mathbb{P}_{p^*}[X \in \mathcal{B}(0, \epsilon)]}\]

\begin{restatable}[Projection Norm Bound]{lemma}{} \label{lemma: pnorm bound} Suppose $X\sim p_0$ and $p^*$ is some spherically symmetric distribution with covariance $\sigma^2I_d$. Let
\[w(X) = \mathbbm{1}\{X \in \mathcal{B}(0,\epsilon)\}\frac{p^*(X)}{p_0(X)}.\]
Then we have
\begin{equation}
   \left\|\left\{\mathbb{E}_{X\sim p_{0}}[w(X)XX^{T}]\right\}^{-1}\mathbb{E}_{X\sim p_{0}}[w(X)X]\right\|_2 \lesssim \epsilon^{-1}.
\end{equation}
\end{restatable}

\noindent \textit{Proof}. Let $b:= \Big(a^{T}\left\{\mathbb{E}_{X\sim p_{0}}[w(X)XX^{T}]\right\}^{-1}\Big)^{T}$. We have,
\begin{align}
    \left\|\left\{\mathbb{E}_{X\sim p_0}[w(X)XX^{T}]\right\}^{-1}\mathbb{E}_{X\sim p_0}[w(X)X]\right\|_2 &= \sup_{\|a\|_2\neq 0} \frac{a^{T}\left\{\mathbb{E}[w(X)XX^{T}]\right\}^{-1}\mathbb{E}[w(X)X]}{\sqrt{a^{T}a}} \\
    &= \sup_{\|b\|_2 \neq 0} \frac{b^{T}\mathbb{E}[w(X)X]}{\sqrt{b^{T} \mathbb{E}\big[w(X)XX^{T}\big]\mathbb{E}\big[w(X)XX^{T}\big]^{T}b}}. \label{eq: projection norm}
\end{align}
First we'll analyze the numerator. Observe that for a any $b$ we can bound it as
\begin{align}
    b^{T}\mathbb{E}_{X \sim p_{0}}[w(X)X] &= b^{T}\mathbb{E}_{X \sim p_{0}}\left[\mathbbm{1}\Big\{X \in \mathcal{B}(0, \epsilon)\Big\}\frac{p^*(X)}{p_{0}(X)}X\right] \\
    &= b^{T}\mathbb{E}_{X \sim p^*}\left[\mathbbm{1}\Big\{X \in \mathcal{B}(0, \epsilon)\Big\}X\right] \\
    &= b^{T}\int_{\mathbb{R}^d}p^*(X)\mathbbm{1}\Big\{X \in \mathcal{B}(0, \epsilon)\Big\}X dx_1 \dots dx_d \\
    &= b^{T}\int_{\mathcal{B}(0, \epsilon)}p^*(X)X dx_1 \dots dx_d \\
    &= \int_{\mathcal{B}(0, \epsilon)}p^*(X)(b^{T}X) dx_1 \dots dx_d \\
    &\leq \epsilon \|b\|_2 \int_{\mathcal{B}(0, \epsilon)}p^*(X) dx_1 \dots dx_d \\
    &= \epsilon\|b\|_2\mathbb{P}_{p^*}(X \in \mathcal{B}(0, \epsilon)). \label{eq: numerator bound}
\end{align}
Similarly, for a any $b$ the denominator can be bounded as
\begin{align}
    \sqrt{b^{T} \mathbb{E}\big[w(X)XX^{T}\big]\mathbb{E}\big[w(X)XX^{T}\big]^{T}b} &= \left\| \mathbb{E}_{X \sim p_0}\big[w(X)XX^{T}\big]^{T}b\right\|_2 \\
    &\geq \sigma_{\min}\left(\mathbb{E}_{X\sim p_0}\big[w(X)XX^{T}\big]\right) \|b\|_2 \\
    &= \sigma_{\min}\left(\mathbb{E}_{X\sim p^*}\left[\mathbbm{1}\Big\{X \in \mathcal{B}(0, \epsilon)\Big\}XX^{T}\right]\right) \|b\|_2. \label{eq: denominator bound 1}
\end{align}
Now it suffices to find a lower bound on the smallest singular value of the restricted scatter matrix. Note that since the matrix is symmetric and positive semidefinite, we have 
\begin{align*}
    \sigma_{\min}\left(\mathbb{E}_{X\sim p^*}\left[\mathbbm{1}\Big\{X \in \mathcal{B}(0, \epsilon)\Big\}XX^{T}\right]\right) &= \lambda_{\min}\left(\mathbb{E}_{X\sim p^*}\left[\mathbbm{1}\Big\{X \in \mathcal{B}(0, \epsilon)\Big\}XX^{T}\right]\right) \\
    &= \min_{v \in \mathbb{S}^{d-1}}v^{T}\mathbb{E}_{X\sim p^*}\left[\mathbbm{1}\Big\{X \in \mathcal{B}(0, \epsilon)\Big\}XX^{T}\right]v \\
    &= \min_{v \in \mathbb{S}^{d-1}}\mathbb{E}_{X\sim p^*}\left[\mathbbm{1}\Big\{X \in \mathcal{B}(0, \epsilon)\Big\}v^{T}XX^{T}v\right] \\
    &= \min_{v \in \mathbb{S}^{d-1}}\mathbb{E}_{X\sim p^*}\left[\mathbbm{1}\Big\{X \in \mathcal{B}(0, \epsilon)\Big\}(X^{T}v)^2\right] \\
    &= \min_{v \in \mathbb{S}^{d-1}} \int_{\mathcal{B}(0, \epsilon)}(X^{T}v)^2p^*(X)dx_1\dots dx_d \\
    & \geq \underbrace{\inf_{Z\in \mathcal{B}(0, \epsilon)}p^*(Z)}_{:=m}\min_{v \in \mathbb{S}^{d-1}} \int_{\mathcal{B}(0, \epsilon)}(X^{T}v)^2dx_1\dots dx_d.
\end{align*}
Now we have
\begin{align}
    \sigma_{\min}\left(\mathbb{E}_{X\sim p^*}\left[\mathbbm{1}\Big\{X \in \mathcal{B}(0, \epsilon)\Big\}XX^{T}\right]\right) &\geq m\min_{v \in \mathbb{S}^{d-1}} \int_{\mathcal{B}(0, \epsilon)}(X^{T}v)^2dx_1\dots dx_d \\
    &= m \min_{v \in \mathbb{S}^{d-1}} v^{T}\left(\int_{\mathcal{B}(0,\epsilon)}XX^{T} dx_1\dots dx_d\right)v \\
    &= m \mu\left(\mathcal{B}(0, \epsilon)\right)  \min_{v \in \mathbb{S}^{d-1}} v^{T}\Sigma_{\mathcal{B}(0,\epsilon)} v
\end{align}
where $\Sigma_{\mathcal{B}(0,\epsilon)}$ is the covariance of a uniform probability measure over $\mathcal{B}(0,\epsilon)$, and $\mu$ denotes the Lebesgue measure on $\mathbb{R}^d$. By \Cref{prop: spherical cov} we have $\Sigma_{\mathcal{B}(0,\epsilon)} = \frac{\epsilon^2}{d+2}I_d$, thus
\begin{align}
    \sigma_{\min}\left(\mathbb{E}_{X\sim p^*}\left[\mathbbm{1}\Big\{X \in \mathcal{B}(0, \epsilon)\Big\}XX^{T}\right]\right) &\geq m \mu\left(\mathcal{B}(0, \epsilon)\right) \epsilon^2 \frac{1}{d+2} \min_{v \in \mathbb{S}^{d-1}} v^{T}I_d v \\
    &= \mu\left(\mathcal{B}(0, \epsilon)\right) \frac{m\epsilon^2}{d+2} \label{eq: denominator bound 2}
\end{align}
Plugging this, \cref{eq: denominator bound 1} and \cref{eq: numerator bound} back into \cref{eq: projection norm}, we obtain
\begin{align*}
    \left\|\left\{\mathbb{E}_{X\sim p_0}[w(X)XX^{T}]\right\}^{-1}\mathbb{E}_{X \sim p_0}[w(X)X]\right\|_2 &\leq \frac{m\epsilon\|b\|_2\mathbb{P}_{p^*}(X \in \mathcal{B}(0, \epsilon))}{\|b\|_2\mu\left(\mathcal{B}(0, \epsilon)\right) \frac{\epsilon^2}{d+2}} \\
    &= \frac{m\mathbb{P}_{p^*}(X \in \mathcal{B}(0, \epsilon))}{\mu\left(\mathcal{B}(0, \epsilon)\right) \frac{\epsilon}{d+2}} \\
    &\leq \frac{m\sup_{X\in \mathbb{R}^d}p^*(X) \mu\left(\mathcal{B}(0, \epsilon)\right)}{\mu\left(\mathcal{B}(0, \epsilon)\right) \frac{\epsilon }{d+2}} \\
    &= \frac{m(d+2)\sup_{X\in \mathbb{R}^d}p^*(X)}{ \epsilon } \\
    &\lesssim \epsilon^{-1}.
\end{align*}
\qed{}
% \begin{align*} 
% \Big\|\theta^*_{x_0, \epsilon} - \tilde{\theta}_{x_0, \epsilon}^*\Big\|_{\Sigma_{x_0}^{\epsilon}} &= \left\|\mathbb{E}[w(X)XX^{T}]^{-1}\mathbb{E}[w(X)X\xi(X)]\right\|_{\Sigma_{x_0}^{\epsilon}} \\
% &\leq \frac{1}{2}\left\|\mathbb{E}[w(X)XX^{T}]^{-1}\mathbb{E}[w(X)X]\right\|_{\Sigma_{x_0}^{\epsilon}}L_fM\epsilon^2 \\
% &= \frac{1}{2}\Big\|\mathbb{E}_{X\sim p^*}\left[\mathbbm{1}\{X \in \mathcal{B}(x_0, \epsilon)\}XX^{T}\right]^{-1}\mathbb{E}_{X\sim p^*}[\mathbbm{1}\{X \in \mathcal{B}(x_0, \epsilon)\}X]\Big\|_{\Sigma_{x_0}^{\epsilon}}L_fM\epsilon^2 \\
% &\leq \frac{1}{2}\cdot\frac{\sigma_{\max}(\Sigma_{x_0}^{\epsilon})}{\sigma_{\min}(\Sigma^{\epsilon}_{x_0})} \Big\|\mathbb{E}_{X\sim p^*}[\mathbbm{1}\{X \in \mathcal{B}(x_0, \epsilon)\}X]\Big\|_2L_fM\epsilon^2 \\
% &\leq \frac{1}{2}\cdot\frac{\sigma_{\max}(\Sigma_{x_0}^{\epsilon})}{\sigma_{\min}(\Sigma^{\epsilon}_{x_0})}BL_fM\epsilon^2 \\
% &\lesssim \epsilon^2.
% \end{align*}



\begin{restatable}[Surrogate]{lemma}{} \label{lemma: surrogate problem} Let 
\[\tilde{\theta}^*_{x_0, \epsilon} = \arg\min_{\theta, \alpha}\mathbb{E}_{X\sim p_{0}}\Big[\mathbbm{1}\Big\{X \in \mathcal{B}(0, \epsilon)\Big\}\frac{p^*(X)}{p_{0}(X)}\Big(g(X + x_0) - X^{T}\theta - \alpha \Big)^2\Big]\]
with 
\[g(X) = f\Big(\Pi(x_0) + \nabla\Pi(x_0)^{T}(X-x_0)\Big).\]
Then we have $\tilde{\theta}^*_{x_0, \epsilon} \propto \nabla \Pi(x_0)$. 
\end{restatable}
\noindent \textit{Proof.} Define 
\[R_{\epsilon}(\theta; X) := \mathbb{E}_{X\sim p_{0}}\Big[\mathbbm{1}\Big\{X \in \mathcal{B}(0, \epsilon)\Big\}\frac{p^*(X)}{p_{0}(X)}\Big(g(X + x_0) - X^{T}\theta - \alpha \Big)^2\Big].\]
Then we have,
\begin{align*}
    R_{\epsilon}(\theta; X)&= \int_{\mathbb{R}^d}\mathbbm{1}\Big\{X \in \mathcal{B}(0, \epsilon)\Big\}\frac{p^*(X)}{p_{0}(X)}\Big(g(X+x_0) - X^{T}\theta - \alpha\Big)^2p_{\epsilon,_0}(X) dz_1 \dots dz_d \\
    &= \int_{\mathbb{R}^d}\mathbbm{1}\Big\{X \in \mathcal{B}(0, \epsilon)\Big\}p^*(X)\Big(g(X + x_0) - X^{T}\theta - \alpha \Big)^2 dz_1 \dots dz_d.
\end{align*}
Now let 
\begin{equation}
    A(\epsilon) := \int_{\mathbb{R}^d} \mathbbm{1}\Big\{X \in \mathcal{B}(0, \epsilon)\Big\}p^*(X) dz_1 \dots dz_d.
\end{equation}
Then we have,
\begin{align*}
    R_{\epsilon}(\theta; X)&= A(\epsilon) \bigintsss_{\mathbb{R}^d}\underbrace{\frac{\mathbbm{1}\Big\{X \in \mathcal{B}(0, \epsilon)\Big\}p^*(X)}{A(\epsilon)}}_{:= \tilde{p}(X)}\Big(g(X+x_0) - X^{T}\theta - \alpha\Big)^2 dz_1 \dots dz_d \\
    &= A(\epsilon)\, \mathbb{E}_{X \sim \tilde{p}}\left(g(X+x_0)-X^{T}\theta - \alpha\right)^2 \\
    &= A(\epsilon)\,\mathbb{E}\left[\mathbb{E}\left[\left(g(X+x_0)-X^{T}\theta - \alpha\right)^2 \, \Big| 
    \, \nabla\Pi(x_0)^{T}X\right]\right]\\
    &= A(\epsilon)\,\mathbb{E}\left[\mathbb{E}\left[\left(f\Big(\Pi(x_0) + \nabla\Pi(x_0)^{T}X\Big)-X^{T}\theta-\alpha\right)^2 \, \Big| 
    \, \nabla\Pi(x_0)^{T}X\right]\right].
\end{align*}
By \Cref{prop: conditional spherical symmetry} we know that $\tilde{p}$ is also a spherically symmetric distribution, and it therefore satisfies the LCE property. Now we'll apply Jensen's inequality to say,
\begin{align*}
    R_{\epsilon}(\theta; X) &\geq A( \epsilon)\,\mathbb{E}\left[\left(f\Big(\Pi(x_0) + \nabla\Pi(x_0)^{T}X\Big)-\mathbb{E}\left[X^{T}\theta \, \Big| 
    \, \nabla\Pi(x_0)^{T}X\right] - \alpha\right)^2\right].
\end{align*}
Since $\tilde{p}$ is spherically symmetric, by eq (2.18) of \cite{shevlyakova2013sliced} we know that 
\begin{align}\mathbb{E}_{\tilde{p}}[X|\nabla\Pi(x_0)^TX] &= \mathbb{E}_{\tilde{p}}[X] + \frac{\Sigma_{\tilde{p}}\nabla\Pi(x_0)\nabla\Pi(x_0)^T(X - \mathbb{E}_{\tilde{p}}[X])}{\nabla\Pi(x_0)^T\Sigma_{\tilde{p}}\nabla\Pi(x_0)} \\
&= \left(\frac{\nabla\Pi(x_0)}{\|\nabla\Pi(x_0)\|_2^2}\right)\nabla\Pi(x_0)^TX
\end{align}
since $\Sigma_{\tilde{p}} \propto I_d$ and $\mathbb{E}_{\tilde{p}}[X] = 0$. Thus,
\begin{align*}
    R_{\epsilon}(\theta; X) &\geq A(\epsilon)\,\mathbb{E}\Bigg[\bigg(f\Big(\Pi(x_0) + \nabla\Pi(x_0)^{T}X\Big)-\underbrace{\left(\frac{\theta^T\nabla\Pi(x_0)}{\|\nabla\Pi(x_0)\|_2^2}\right)}_{:=C(\theta)}\nabla\Pi(x_0)^TX -\alpha\bigg)^2\Bigg] \\
    &= A(\epsilon)\,\mathbb{E}\left[\left(f\Big(\Pi(x_0) + \nabla\Pi(x_0)^{T}X\Big)-C(\theta)\cdot\nabla\Pi(x_0)^{T}X -\alpha\right)^2\right].
\end{align*}
Observe that the inequality becomes an equality when $\theta = C(\theta) \cdot\nabla\Pi(x_0)$, which completes the proof.

\qed{}

\begin{restatable}[Conditional Spherical Symmetry]{prop}{} \label{prop: conditional spherical symmetry}
    Suppose $X$ has a density $f(X)$ satisfying spherical symmetry (\Cref{def: spherical symmetry}) with mean $0$ and covariance $\sigma^2I$ for some $\sigma > 0$. Then for any $\epsilon > 0$ the distribution of the random variable $X \,| \, X \in \mathcal{B}(0, \epsilon)$ is also spherically symmetric. 
\end{restatable}
\noindent \textit{Proof.} We have that for any orthogonal $Q$ and for any Borel set $A$
\begin{align*}
    \mathbb{P}(QX \in A \, | \, X \in \mathcal{B}(0,\epsilon) ) &= \int_{A}f(Qz\,|\,X \in \mathcal{B}(0,\epsilon))dz_1\dots dz_d \\
    &= \bigintsss_{A}\frac{f(Qz) \mathbbm{1}\{z \in \mathcal{B}(0,\epsilon)\}}{\int_{\mathcal{B}(0,\epsilon)}f(z')dz'_1 \dots dz'_d}dz_1\dots dz_d \\
    &= \bigintsss_{A}\frac{f(z) \mathbbm{1}\{z \in \mathcal{B}(0,\epsilon)\}}{\int_{\mathcal{B}(0,\epsilon)}f(z')dz'_1 \dots dz'_d}dz_1\dots dz_d
\end{align*}
where the last line follows from spherical symmetry of $X$. Thus,
\begin{align*}
    \mathbb{P}(QX \in A \, | \, X \in \mathcal{B}(0,\epsilon) ) &= \bigintsss_{A}\frac{f(z) \mathbbm{1}\{z \in \mathcal{B}(0,\epsilon)\}}{\int_{\mathcal{B}(0,\epsilon)}f(z')dz'_1 \dots dz'_d}dz_1\dots dz_d \\
    &= \int_{A}f(z\,|\,X \in \mathcal{B}(0,\epsilon))dz_1\dots dz_d \\
    &= \mathbb{P}(X \in A \, | \, X \in \mathcal{B}(0,\epsilon) ).
\end{align*}
\qed{}

\begin{restatable}[Approximation]{lemma}{} \label{lemma: proj approx}
    Suppose $f$ is $L_f$-Lipschitz. Then for Lebesgue almost every $x_0 \in \mathcal{B}(0,r)$, there exists finite $\epsilon, \lambda > 0$ such that for all $x \in \mathcal{B}(x_0, \epsilon)$
    \[\bigg|f\big(\Pi(x)\big) - f\big(\Pi(x_0) + \nabla\Pi(x_0)^{T}(x-x_0) \big)\bigg| \lesssim (M_2+rM_3)L_f\epsilon^2. \]
\end{restatable}

\noindent \textit{Proof.} By \Cref{lemma: operator norm bound} we have that for Lebesgue almost every $x_0 \in \mathcal{B}(0,r)$, there exists finite $\epsilon, \lambda > 0$ such that for every $x \in \mathcal{B}(x_0, \epsilon)$
\[\|\nabla^2\Pi(x)\|_{\text{op}} \leq \underbrace{2M_2\lambda^{2}+rM_3\lambda^{3}}_{:= C_{x_0, \lambda}}.\]
We can employ an expansion of the projection map as follows,
\begin{align*}
    \Pi(x) &= \Pi(x_0) + \nabla\Pi(x_0)^{T}(x-x_0) + \underbrace{\int_{0}^{1}(x-x_0)^T\nabla^2\Pi(x+t(x-x_0))(x-x_0) dt}_{:= R(x)}.
\end{align*}
By the operator norm bound, we have that
\[\|R(x)\|_2 \leq \int_{0}^{1}\|x - x_0\|_2^2\|\|\nabla^2\Pi(x+t(x-x_0))\|_{\text{op}} \,dt \leq C_{x_0, \lambda}\epsilon^2.\]
It follows that 
\begin{align}
    \left|\,\Pi(x) - \left(\Pi(x_0) + \nabla\Pi(x_0)^{T}(x-x_0)\right)\right| &\leq C_{x_0, \lambda}\epsilon^2.
\end{align}
The lemma follows from the Lipschitz continuity of $f$. 
\qed{}


\begin{restatable}[Perturbation]{lemma}{} \label{lemma: perturbation}
    Let 
    \[\theta^*_{\alpha} = \arg\min_{\theta}\mathbb{E}\big[w(X)\big(Y - \theta^{T}X - \alpha\big)^2\big]\]
    and define $\tilde{\theta}^*$ to be the solution to the same optimization with perturbed responses,
    \[\tilde{\theta}^*_{\alpha} = \arg\min_{\theta}\mathbb{E}\big[w(X)\big(Y - \xi(X) - \theta^{T}X - \alpha\big)^2\big].\]
    If $\mathbb{E}[w(X)XX^T] \succ 0$, then we have
    \[\theta^* - \tilde{\theta}^* = \mathbb{E}[w(X)XX^{T}]^{-1}\mathbb{E}[w(X)X\xi(X)].\]
\end{restatable}
\noindent \textit{Proof.} We can use the normal equations to say
\begin{align*}
    \mathbb{E}[w(X)XX^{T}]\theta^*_{\alpha} = \mathbb{E}[w(X)X(Y-\alpha)]
\end{align*}
implying 
\[\theta^*_{\alpha} = \mathbb{E}[w(X)XX^{T}]^{-1}\mathbb{E}[w(X)X(Y-\alpha)].\]
We can apply the same argument to obtain 
\[\tilde{\theta}^*_{\alpha} = \mathbb{E}[w(X)XX^{T}]^{-1}\mathbb{E}[w(X)X(Y - \xi(X) -\alpha )].\]
Taking the difference, we obtain
\begin{align*}
    \theta^*_{\alpha} - \tilde{\theta}^*_{\alpha} &= \mathbb{E}[w(X)XX^{T}]^{-1}\mathbb{E}[w(X)X(Y - \alpha -  Y + \xi(X) +\alpha )] \\
    &= \mathbb{E}[w(X)XX^{T}]^{-1}\mathbb{E}[w(X)X\xi(X)].
\end{align*}
\qed{}

\begin{restatable}[LCE Proportionality, \cite{li1989regression}]{thm}{} \label{thm: lce proportionality}
    Suppose $X$'s distribution satisfies the LCE property, $\mathbb{E}[Y|X] = f(b^{T}X)$ for some $b \in \mathbb{R}^d$ and 
    \[\theta^*, \alpha^* = \arg\min_{\theta, \alpha}\mathbb{E}\big(f(b^{T}X) - \theta^{T}X - \alpha\big)^2.\]
    Then $\theta^* \propto b$. 
\end{restatable}
\noindent \textit{Proof}. We have 
\begin{align*}
    \mathbb{E}\big(f(b^{T}X) - \theta^{T}X-\alpha\big)^2 &= \mathbb{E}\Big[ \mathbb{E}\Big[\big(f(b^{T}X) - \theta^{T}X-\alpha\big)^2 \, \Big | \, b^{T}X\Big] \Big] \\
    &\geq \mathbb{E}\Big[ \Big(f(b^{T}X) - \mathbb{E}\big[\theta^{T}X |b^{T}X\big]-\alpha\Big)^2 \Big] \\
    &= \mathbb{E}\Big[ \Big(f(b^{T}X) - Cb^{T}X-\alpha'\Big)^2 \Big]
\end{align*}
where the penultimate line uses Jensen's inequality and the last line uses the LCE property. Observe that the inequality becomes an inequality when $\theta = Cb$ and $\alpha = \alpha'$ from which the theorem follows. 

\qed{}

\begin{restatable}[Projection Gradient]{prop}{} \label{prop: proj gradient}
    Let $\Pi: \mathbb{R}^d \rightarrow [0, \text{len}_{\gamma}]$ be the closest-point projection onto $\gamma$ as defined in \Cref{def: projection}. Then for every $x \notin \overline{\text{Sing}(\Gamma)}$ such that 
    \[\langle\gamma''(\Pi(x)), x - \gamma(\Pi(x))\rangle \neq 1\]
    we have 
    \[\nabla\Pi(x)  = -\frac{\gamma'(\Pi(x))}{-1 + \left\langle\gamma''(\Pi(x)), x - \gamma(\Pi(x))\right\rangle}.\]
\end{restatable}

\noindent \textit{Proof of \Cref{prop: proj gradient}.} Let $f(x,t) := \|x - \gamma(t)\|_2^2$. By non-negativity of $\|\cdot\|_2$ and monotonicity of $h(x) = x^2$ for $x \geq 0$, we have $\Pi(x) = \arg\min_t \|x - \gamma(t)\|_2 = \arg \min_t \|x - \gamma(t)\|_2^2 = \arg\min_t f(x, t)$. We have the first-order condition that 
\[ \partial_t f(x,t) = 0\]
at the optimum. Now, if we define $F(x, t) := \partial_t f(x,t)$ the first order condition amounts to the requirement that $F(x,t) = 0$. Expanding $F$, we have
\begin{align}
    F(x, t) &= \partial_t \|x - \gamma(t)\|_2^2 \\ \label{eq: ift}
    &= -2\langle x - \gamma(t)), \gamma'(t)\rangle.
\end{align}
Observe that the first order condition implies that the projection displacement $x - \gamma(t)$ lies normal to the curve at optimality. 
Now we'll apply the implicit function theorem to $F(x, t)$. Because $x$ is \textit{not} in the closure of the singular set, we have that $\Pi(x)$ is unique and $\Pi$ is differentiable in a neighborhood about $x$ - its not hard to verify that this yields continuous differentiability of $F(x,t)$ in a neighborhood about $(x, \Pi(x))$. This, coupled with the assumption that $\langle\gamma''(\Pi(x)), x - \gamma(\Pi(x))\rangle \neq 1$ allows us apply the implicit function theorem to conclude that 
\begin{align}
    \nabla\Pi(x) &= - \frac{\nabla_x F\big(x, t\big)}{\partial_t F\big(x, t\big)} \Bigg|_{t = \Pi(x)}\\
    &= - \frac{\nabla_x\left(-2\langle x - \gamma(t), \gamma'(t)\rangle\right)}{\partial_t \left(-2\langle x - \gamma(t), \gamma'(t)\rangle\right)} \Bigg|_{t = \Pi(x)}\\
    &= -\frac{\gamma'(t)}{-\|\gamma'(t)\|_2^2 + 
    \langle\gamma''(t), x - \gamma(t)\rangle}\Bigg|_{t = \Pi(x)} \label{eq: projection grad} \\ 
    &= -\frac{\gamma'(\Pi(x))}{-1 + \left\langle\gamma''(\Pi(x)), x - \gamma(\Pi(x))\right\rangle}.
\end{align}
\qed{}


\begin{restatable}[Projection Hessian]{prop}{} \label{prop: proj hessian}
    Let $\Pi: \mathbb{R}^d \rightarrow [0, \text{len}_{\gamma}]$ be the closest-point projection onto $\gamma$ as defined in \Cref{def: projection}. Then for every $x \notin \overline{\text{Sing}(\Gamma)}$ and 
    \[\langle\gamma''(s), x - \gamma(s)\rangle \neq 1\]
    we have
    \[\nabla^2\Pi(x) = \frac{1}{g^2(x)}\left(\gamma''(s)\gamma'(s)^T + \gamma'(s)\gamma''(s)^T\right)- \frac{\langle\gamma^{(3)}(s), x - \gamma(s)\rangle}{g^3(x)}\gamma'(s)\gamma'(s)^T.\]
    where $s = \Pi(x)$ and $g(x) = -1 + \langle\gamma''(s), x-\gamma(s)\rangle$.
\end{restatable}
\noindent \textit{Proof.} By \cref{prop: proj gradient} we know that for all $x$ such that the assumptions hold, we have
\begin{align}
    \nabla\Pi(x) &= -\frac{\gamma'(s)}{-1 + \langle\gamma''(s), x - \gamma(s)\rangle} \\
    &= -\frac{\gamma'(s)}{g(x)}.
\end{align}
Taking the gradient again yields
\begin{align}
    \nabla^2\Pi(x) &= -\frac{g(x)\nabla \gamma'(s) - \gamma'(s)\nabla g(x)^T}{g(x)^2} \\
    &= -\frac{\nabla \gamma'(s)}{g(x)} + \frac{1}{g^2(x)}\gamma'(s)\nabla g(x)^T.
\end{align}
By the chain rule, we have $\nabla\gamma'(s) = \gamma''(s)\nabla\Pi(x)^T$ and 
\begin{align}
    \nabla g(x) &= \nabla\langle\gamma''(s), x\rangle -\nabla\langle \gamma''(s), \gamma(s)\rangle \\
    &= \langle \gamma^{(3)}(s), x\rangle \nabla\Pi(x) + \gamma''(s) - \langle\gamma^{(3)}(s), \gamma(s)\rangle\nabla\Pi(x) - \underbrace{\langle\gamma''(s), \gamma'(s)\rangle}_{= 0}\nabla\Pi(x) \\
    &= \langle\gamma^{(3)}(s), x - \gamma(s)\rangle\nabla\Pi(x) + \gamma''(s).
\end{align}
Thus,
\begin{align*}
    \nabla^2\Pi(x) &= -\frac{1}{g(x)}\gamma''(s)\nabla\Pi(x)^T + \frac{\langle\gamma^{(3)}(s), x - \gamma(s)\rangle}{g^2(x)}\gamma'(s)\nabla\Pi(x)^T + \frac{1}{g^2(x)}\gamma'(s)\gamma''(s)^T \\
    &= \frac{1}{g^2(x)}\gamma''(s)\gamma'(s)^T - \frac{\langle\gamma^{(3)}(s), x - \gamma(s)\rangle}{g^3(x)}\gamma'(s)\gamma'(s)^T + \frac{1}{g^2(x)}\gamma'(s)\gamma''(s)^T \\
    &= \frac{1}{g^2(x)}\left(\gamma''(s)\gamma'(s)^T + \gamma'(s)\gamma''(s)^T\right)- \frac{\langle\gamma^{(3)}(s), x - \gamma(s)\rangle}{g^3(x)}\gamma'(s)\gamma'(s)^T.
\end{align*}
\qed{}

\begin{restatable}[Operator Norm Bound]{lemma}{} \label{lemma: operator norm bound}
    Suppose $x \in \mathcal{B}(0,r) \cap R$ with $R$ defined in \cref{eq: good set}. Then there exist finite $\epsilon, \lambda > 0$ such that for every $y\in \mathcal{B}(x,\epsilon)$ 
    \[\|\nabla^2\Pi(y)\|_{\text{op}} \leq 2M_2\lambda^{2}+rM_3\lambda^{3}.\]
\end{restatable}

\noindent \textit{Proof.} By \Cref{thm: differentiability of pi} we have that for every $x \in R$, 
\[\nabla^2\Pi(x) = \frac{1}{g^2(x)}\left(\gamma''(s)\gamma'(s)^T + \gamma'(s)\gamma''(s)^T\right)- \frac{\langle\gamma^{(3)}(s), x - \gamma(s)\rangle}{g^3(x)}\gamma'(s)\gamma'(s)^T.\]
where $s = \Pi(x)$ and $g(x) = -1 + \langle\gamma''(s), x-\gamma(s)\rangle$. Since $R$ is an open set, there exists an $\epsilon > 0$ such that this expression holds for all $y \in \mathcal{B}(x, \epsilon)$.  Bounding the operator norm yields
\begin{align*}
    \|\nabla^2\Pi(x)\|_{\text{op}} &= \max_{v\in \mathbb{S}^{d-1}} \|\nabla^2\Pi(x)v\|_2 \\
    &\leq \max_{v\in \mathbb{S}^{d-1}}\left(\frac{\|\gamma''(s)\gamma'(s)^Tv + \gamma'(s)\gamma''(s)^Tv\|_2}{|g^2(x)|} + \frac{\langle\gamma^{(3)}(s), x - \gamma(s)\rangle}{|g^3(x)|} \|\gamma'(s)\gamma'(s)^Tv\|_2\right) \\
    &\leq \max_{v\in \mathbb{S}^{d-1}}\left(\frac{\|\gamma''(s)\gamma'(s)^Tv\|_2 + \|\gamma'(s)\gamma''(s)^Tv\|_2}{|g^2(x)|} + \frac{\langle\gamma^{(3)}(s), x - \gamma(s)\rangle}{|g^3(x)|} \|\gamma'(s)\gamma'(s)^Tv\|_2\right) \\
    &\leq \frac{2M_2}{|g^2(x)|} + \frac{\langle\gamma^{(3)}(s), x - \gamma(s)\rangle}{|g^3(x)|} \\
    &\leq \frac{2M_2}{|g^2(x)|} + \frac{r M_3}{|g^3(x)|}.
\end{align*}
where the final line follows from the assumption that $\Gamma = \text{Im}(\gamma) \subset \mathcal{B}(0, r)$ and $x \in \mathcal{B}(0, r)$. Now we will establish that $1/|g(x)|$ is continuous everywhere in $\mathcal{B}(0,r) \cap R$. This property is a consequence of 
\begin{enumerate}
    \item the of $|g(x)|$ in $\mathcal{B}(0,r) \cap R$, which follows from the definitions of $g$ and $R$
    \item the fact that $|g(x)| \neq 0$ for all $x \in R$ -- this too follows from the definitions of $g$ and $R$.
\end{enumerate}

From these two properties, we can deduce that $1/|g(x)|$ is continuous everywhere in $\mathcal{B}(0,r) \cap R$. The continuity of $1/|g(x)|^k$ for $k\geq 1$ follows. We will leverage what has been shown to say that for every $x$ in $\mathcal{B}(0,r) \cap R$, $\alpha := 1/|g(x)| <\infty$. Furthermore, for the $\epsilon$ defined earlier, there exists a $\delta>0$ such that for all $y \in \mathcal{B}(x, \epsilon)$, $1/|g(y)| \leq \alpha+\delta := \lambda.$ The statement follows.

\qed{}

\begin{restatable}[IFT Degeneracy]{lemma}{} \label{lemma: IFT degen}
    Let 
    \begin{equation} 
    B:= \{x \in \mathbb{R}^d : 1 - \langle\gamma''(t), x-\gamma(t)\rangle = 0 \text{ for all }t \in \arg\inf_t \|x-\gamma(t)\|_2\} \label{eq: ift degeneracy set}
    \end{equation}
    and let $\mu$ denote the $d$-dimensional Lebesgue measure on $\mathbb{R}^d$. Then $\mu(B) = 0.$
\end{restatable}
\noindent \textit{Proof of \Cref{lemma: IFT degen}.} Let $N\gamma := \{(p, v) : p \in \gamma([0, \text{len}_{\gamma}]), v \in (T_p \gamma)^{\bot}\}$ denote the normal bundle of the curve $\gamma$, where $T_p\gamma$ denotes the tangent space of $\gamma$ at point $p$. Observe that $N\gamma$ forms a smooth $d$-dimensional manifold - explicitly, one has the canonical local chart $(t,w) \rightarrow (\gamma(t), w_1n_1(t) + \dots+w_{d-1}n_{d-1}(t))$ where $\{n_i(t)\}_{i\in [d-1]}$ form an orthonormal basis of $N_p\gamma$. 

Now let $B':= \{(p,v) \in N\gamma : \langle v,  \gamma''(\gamma^{-1}(p))\rangle = 1\}$, and let $\Phi:N\gamma \rightarrow \mathbb{R}$ be the function
\[\Phi(p,v) = \langle v, \gamma''(\gamma^{-1}(p))\rangle - 1.\]
Observe that for all $x' \in B$, for all $t' \in \arg\min_t \|x-\gamma(t)\|$ we have \[\langle \underbrace{x' - \gamma(t')}_{v'}, \gamma''(t')\rangle = 1\] where $v' \in N_{\gamma(t')}\gamma$. Thus $x' \in \exp(B')$ where $\exp(p,v) = p + v$. It follows that $B\subset \exp(B')$.

We will now show that $0$ is a regular value of $\Phi$, allowing us to conclude that the level set $\Phi^{-1}(0) = B$ is an embedded submanifold of $N\gamma$ of codimension 1. To show that $0$ is a regular value of $\Phi$, observe that $\Phi$ can be expressed in local coordinates as
\[\Phi(t,w) = \sum_{i = 1}^{d-1}w_i\langle n_i(t), \gamma''(t)\rangle - 1.\]
Taking the partial derivative with respect to $w_i$ we have
\[\partial_{w_i}\Phi(t,w) = \langle n_i(t), \gamma''(t)\rangle\]
To show the regularity of $\Phi^{-1}(0)$, we need to establish that for all $(t,w)$ such that $\Phi(t,w) = 0$, the gradient is nonzero. Observe that for all such $(t,w)$, $\langle v(t,w), \gamma''(t)\rangle = 1$, where $v(t,w) = \sum_{i = 1}^{d-1}w_in_{i}(t)$ -- therefore $\gamma''(t) \neq 0$. Now we need to guarantee that $\gamma''(t)$ has a component in $N_{\gamma(t)}\gamma$. We will show this by contradiction: suppose $\gamma''(t) \in T_{\gamma(t)}\gamma$. Then $\langle\gamma''(t), v\rangle = 0$ for all $v$ normal to $\gamma$ at $t$, which contradicts $\langle \gamma''(t), v(t,w) \rangle = 1$.

Since $\gamma''(t)$ necessarily has a some component normal to the curve, we know that $\partial_{w_i}\Phi(t,w)$ is nonzero for some $i$. This renders the differential of $\Phi$ surjective for all $(t,w)$ such that $\Phi(t,w) = 0$. This means one can apply the regular level set theorem to conclude that $B'$ forms a $d-1$ dimensional submanifold of the normal bundle of $\gamma$ (\citet{lee2003smooth}, Corollary 5.14). Now since $\exp\big|_{B'}$ is a smooth map from $B'$ to $\mathbb{R}^d$, $\exp(B')$ must have $d$-dimensional Lebesgue measure $0$ (\citet{lee2003smooth} Corollary 6.11). Since $B \subset \exp(B')$, by monotonicity of the Lebesgue measure $B$ must have measure $0$ as well.

\qed{}


\begin{restatable}[IFT Degeneracy is Closed]{lemma}{} \label{lemma: IFT degen closed}
    The set 
    \begin{equation}
    B = \{x \in \mathbb{R}^d : 1 - \langle\gamma''(t), x-\gamma(t)\rangle = 0 \text{ for all }t \in \arg\inf_t\|x-\gamma(t)\|_2\}
    \end{equation}
    is a closed subset of $\mathbb{R}^d$. 
\end{restatable}

\noindent \textit{Proof.} Define the set
\begin{equation}
S:= \{(x,t) \in \mathbb{R}^d\times I : t \in \arg\inf_{t\in I}\|x-\gamma(t)\|_2, \, 1 - \langle\gamma''(t), x - \gamma(t)\rangle = 0\} \label{eq: S}
\end{equation}
which is trivially a subset of $\mathbb{R}^{d+1}$. If one fixes $s \in I$, the map $(x,t) \mapsto \|x-\gamma(t)\|_2 - \|x-\gamma(s)\|_2$ is continuous. Since sublevel sets of continuous functions are closed, the set
\begin{equation}
S_s = \{(x,t) : \|x-\gamma(t)\|_2 - \|x-\gamma(s)\|_2 \leq 0\}
\end{equation}
is a closed subset of $\mathbb{R}^{d+1}$. Observe that
\begin{align*}
    \bigcap_{s \in I} S_s &= \bigcap_{s \in I} \Big\{(x,t) : \|x-\gamma(t)\|_2 - \|x-\gamma(s)\|_2 \leq 0\Big \} \\
    &= \Big\{(x,t) : t\in \arg\inf_{t\in I}\|x-\gamma(t)\|_2 \Big \} \\
    &= S'.
\end{align*}
To see why consider the following. The element $(x,t)$ is in $S'$ if and only if $\|x - \gamma(t)\|_2  \leq \|x - \gamma(s)\|_2$ for all $s \in I$. Since the image of $\gamma$ is closed, this occurs if and only if $t \in \arg\inf_{t \in I} \|x - \gamma(t)\|_2$. Having established the closure of $S_s$, the closure of $S'$ follows from the fact that uncountable intersections of closed sets are still closed. 

To establish the closure of the set $S$ defined in \cref{eq: S}, consider the map $(x,t) \mapsto 1 - \langle \gamma''(t), x - \gamma(t)\rangle$. Since $\gamma \in C^3(I)$ this is a continuous map. It follows that the level set $\{(x,t) : 1 - \langle \gamma''(t), x - \gamma(t)\rangle = 0\}$ is closed. Taking the intersection with $S'$ implies the closure of the set $S$. 

Now we will use this to show that $B$ is a closed subset of $\mathbb{R}^d$. To see this, observe that $B$ can be rewritten as $B = \{x : (x,t) \in S\}.$ We will show that $B$ is closed by showing that it contains all of its limit points via its connection to the closed set $S$. Let $\tilde{x}$ be a limit point of $B$. This means there exists a sequence $x_n \in B$ such that $x_n\rightarrow \tilde{x}$. By the connection to $S$ established above, we know that there are associated values $t_n$ such that 
\[t_n \in \arg\inf_{t \in I}\|x_n - \gamma(t_n)\|_2 \quad \text{and} \quad 1 - \langle\gamma''(t_n), x_n - \gamma(t_n)\rangle = 0.\]
Thus, we have a sequence $(x_n, t_n) \in S$ such that $x_n \rightarrow x$. Since $I$, the domain of $\gamma$, is compact, we know that there must exist a subsequence $t_{n_k}$ such that $t_{n_k}\rightarrow t$ for some $t \in I$. Since $x_n$ converges to $x$, we know that the subsequence $x_{n_k}$ also converges to $x$. Thus, we have a sequence $(x_{n_k}, t_{n_k}) \rightarrow (x,t)$, which must be in $S$ due to closure. It follows that $x$ must also be in $B$, and $x_{n_k} \rightarrow x$. Since every subsequence of a convergent sequence converges to the same value, we know that the limit point of $x_{n}$ (which we defined earlier as $\tilde{x}$) is equal to $x$, which is in $B$. Thus, $B$ necessarily contains all of its limit points, and it is therefore closed.  

\qed{}


\begin{restatable}[Closure of the singular set has zero Lebesgue measure]{lemma}{} \label{lemma: singular set measure}
    The closure of the set 
    \begin{equation}\text{Sing}(\Gamma) = \{x \in \mathbb{R}^d: \Pi(x) \text{ fails to be differentiable}\} \label{eq: singular set}
    \end{equation}
    has $d$-dimensional Lebesgue measure $0$. 
\end{restatable}
\noindent \textit{Proof.} Let $d(x, \Gamma)$ be the distance from $x$ to the set $\Gamma$. We will first show that the closure of the set where $d^2(x,\Gamma)$ fails to be differentiable has measure zero. Call this set $\text{Sing}'(\Gamma)$. To show this, observe that the image of the curve $\Gamma$ forms a 1-dimensional $C^3$ submanifold of $\mathbb{R}^d$ with boundary
\[\partial\Gamma = \{\gamma(0), \gamma(\text{len}_{\gamma})\}.\]
Since the boundary consists of a union of two disjoint points, it forms a smooth $0$-dimensional submanifold of $\mathbb{R^d}$. This allows us to apply Remark 4.2 and Corollary 4.12 of \citet{mantegazza2002hamilton} to conclude that the Hausdorff dimension of $\overline{\text{Sing}'(\Gamma)}$ is at most $d-1$. This implies that its $d$-dimensional Hausdorff measure is $0$, which in turn implies that its $d$-dimensional Lebesgue outer measure is $0$. Since the set $\overline{\text{Sing}'(\Gamma)}$ is closed, it is necessarily Borel. It follows that the set is Lebesgue measurable with measure $0$. 

Now we will show that $\Pi(x)$ failing to be differentiable implies $d^2(x, \Gamma)$ fails to be differentiable. To see this observe that
\begin{align}
    d^2(x, \Gamma) &= \inf_{p \in \Gamma}\|x - p\|_2^2 \\ \label{eq: distance to Gamma}
    &= \|x - \gamma(t')\|_2^2
\end{align}
for any $t' \in \Pi(x)$. If $\Pi(x)$ is a singleton, then $d^2(x,\Gamma) = \|x - \gamma(\Pi(x))\|_2^2$ and we have that $d^2(x,\Gamma)$ is not differentiable at $x$ if $\Pi$ is not differentiable at $x$. Now suppose $\Pi(x)$ is not a singleton. We will show that this implies $d^2(x,\Gamma)$ is not differentiable by contradiction. When $\Pi(x)$ is not a singleton then $x$ has at least two distinct nearest points -- pick any two and call them $p$ and $q$. If $d^2(x,\Gamma)$ were differentiable at $x$ then its gradient would be unique and equal to $(x-p')$ where $p'$ realizes the infimum above. Since $p$ and $q$ both realize the infimum, the uniqueness of the gradient of $d^2(x, \Gamma)$ implies that $(x - p) = (x-q)$, which is a contradiction.

This establishes that if $\Pi$ fails to be differentiable at $x$, then $d^2(x, \Gamma)$ must also fail to be differentiable at $x$. This implies $\text{Sing}(\Gamma) \subseteq \text{Sing}'(\Gamma)$, which in turn implies the closure of the former is contained in the closure of the latter. Earlier we established that $\overline{\text{Sing}'(\Gamma)}$ has measure zero -- the Lemma statement therefore follows from monotonicity of the Lebesgue measure.

\qed{}



% \begin{restatable}[Medial axis is closed]{lemma}{}
%     The medial axis of the curve $\gamma$ is closed. 
% \end{restatable}
% \noindent \textit{Proof.} The medial axis $M(\Gamma)$ of $\Gamma = \text{Im}(\gamma)$ can be expressed as
% \[M(\Gamma) = \{x \in \mathbb{R}^d : |\{y \in \Gamma : \|x-y\|_2 = d(x,\Gamma)\}| > 1\}.\]
% To show that $M(\Gamma)$ is closed, we will show that it contains all of its limit points. Consider any sequence $x_n \in M(\Gamma)$ and suppose $x_n \rightarrow x$. We will show that $x$ is necessarily in $M(\Gamma)$. To do so, observe that $\Gamma \times \Gamma$ is compact as $\Gamma$ is compact. Since each $x_n$ is in $M(\Gamma)$, we know it has at least two closest points; we'll pick any two and call them $q_n$ and $p_n$ - this defines two separate sequences $\{q_n\}_{n \in \mathbb{N}}$ and $\{p_n\}_{n \in \mathbb{N}}$, but it also defines the joint sequence $\{(q_n, p_n)\}_{n\in \mathbb{N}} \subset \Gamma\times \Gamma$. By the compactness of this product space, there must exist a convergent subsequence $(q_{n_k}, p_{n_k}) \rightarrow (q,p) \in \Gamma \times \Gamma$. This, in turn, induces a subsequence $x_{n_k}$, which necessarily converges to $x$ since we assume that $x_n$ is a convergent sequence. Now we'll consider two cases. 
% \begin{enumerate}
%     \item \textbf{Case 1:} $(p \neq q)$. Observe that we have the equalities
%     \[d(x_{n_k}, \Gamma) = \|x_{n_k} - p_{n_k}\|_2 = \|x_{n_k} - q_{n_k}\|_2\]
%     by definition. By continuity, this implies that as $k \rightarrow \infty$,
%     \[d(x,\Gamma) = \|x - p\|_2 = \|x - q\|_2\]
%     with $p \neq q$. This directly implies that $x \in M(\Gamma)$. 
%     \item \textbf{Case 2:} $(p = q)$. We will show that, under the hypothesis of positive reach, this cannot happen. Since $x_{n_k} \in M(\Gamma)$, we know that $d(x_{n_k}, \Gamma) \geq \tau_{\gamma}$ for all $k$. This implies that $d(x, \Gamma) \geq \tau_{\gamma}$ as well. 
%     \begin{enumerate}
%         \item $[]$
%     \end{enumerate}
% \end{enumerate}


% Let $f:\mathbb{R}^d\times I\rightarrow \mathbb{R}$ be defined as $f(x,t) = \|x-\gamma(t)\|_2^2$. Then the medial axis $M$ can be expressed as
% \[M = \{x \in \mathbb{R}^d : |\arg\min_{t}f(x,t)| > 1\}.\]
% To show that $M$ is closed, consider the set 
% \[P := \{(x,t,s) \in \mathbb{R}^d \times I \times I: t\neq s, f(x,t) = f(x,s) \leq f(x,u) \, \forall \, u \in I\}.\]
% Observe that $(x,t,s) \in P$ if and only if $t$ and $s$ are distinct parameters of minimal distance points to $x$. Now consider the set
% \[U_{1,u} := \{(x,t,s): f(x,t) - f(x,u) \leq 0\}.\]
% For a fixed $u$, the map $(x,t,s) \mapsto f(x,t) - f(x,u)$ is continuous since $\gamma \in C^3(I)$, rendering the sublevel set closed. Now observe that
% \begin{align}
% U_1 &:= \bigcap_{u \in I}U_{1,u} \\
% &= \left\{(x,t,s) : f(x,t) \leq f(x,u) \, \forall \, u \in I\right\}
% \end{align}
% is closed, since uncountable intersections of closed sets are closed. By a similar argument, one can show that the set 
% \[U_2 := \{(x,t,s): f(x,s) \leq f(x,u) \, \forall \, u \in I \}\]
% is also closed. Finally, define 
% \[U_3 := \{(x,t,s): f(x,t) = f(x,s)\}\]
% which is also closed, as it consists of the zero level set of a continuous function. 

% \qed{}

% \begin{restatable}[Continuous differentiability of $\Pi$]{lemma}{}
%     For any $x$ such that $\Pi(x)$ is differentiable in a neighborhood about $x$, $\Pi(x)$ must also be \textit{continuously} differentiable on that neighborhood.
% \end{restatable}
% \noindent \textit{Proof.} Let $d^2(x,\Gamma)$ be defined as in \cref{eq: distance to Gamma}. Under the hypothesis that $\Pi(x)$ is differentiable in a neighborhood $U$ about $x$, we know that it must also be unique and well defined on $U$. It follows that
% \begin{align}
%     d^2(x,\Gamma) = \|x - \gamma(\Pi(x))\|_2^2 \text{ for all }x \in U.
% \end{align}
% The hypothesis of the lemma coupled with the fact that $\gamma \in C^3(I)$ implies that $d^2(x,\Gamma)$ is differentiable on $U$. 


\begin{restatable}[Covariance of $\text{Unif}(\mathcal{B}(0,\epsilon))$]{prop}{} \label{prop: spherical cov}
    Suppose we have $X$ distributed according to a uniform probability measure over $\mathcal{B}(0,\epsilon)$ in $\mathbb{R}^d$. Then the covariance matrix of $X$ is $\frac{\epsilon^2}{d+2}I_d$. 
\end{restatable}
\noindent \textit{Proof.} By symmetry, the covariance matrix must be a multiple of $I_d$. To find the constant, we need to understand the distribution of $\mathbb{E}\|X\|_2$ - thus we will compute the radial pdf. Note that the cdf of $\|X\|_2$ can be written as
\begin{align*}
    \mathbb{P}(\|X\|_2 \leq r) &= \frac{\mu(\mathcal{B}(0,r))}{\mu(\mathcal{B}(0,\epsilon))} \\
    &= \left(\frac{r}{\epsilon }\right)^d
\end{align*}
where $\mu$ denotes the $d$-dimensional Lebesgue measure on $\mathbb{R}^d$. The radial pdf can be obtained by differentiating the cdf with respect to $r$,
\begin{align*}
    f_r(r) &= \frac{d}{dr}\mathbb{P}(\|X\|_2\leq r) \\
    &= \frac{dr^{d-1}}{\epsilon^d}.
\end{align*}
Now we can compute the covariance scaling factor as 
\begin{align*}
    \mathbb{E}[X_i^2] &= \frac{1}{d}\mathbb{E}\|X\|_2^2 \\
    &=\frac{1}{\epsilon^d}\int_{0}^{\epsilon}r^2r^{d-1}dr \\
    &= \frac{1}{\epsilon^d}\frac{1}{d+2}r^{d+2} \Bigg|^{r=\epsilon}_{r=0} \\
    &= \frac{\epsilon^2}{d+2}.
\end{align*}
\qed{}


\bibliography{main}

\end{document}
